\documentclass[a3paper,landscape]{article}
\usepackage[top=1cm,bottom=1cm,left=1cm,right=1cm]{geometry}
\usepackage[T1]{fontenc}
\usepackage{lmodern}
\usepackage{forest}
\usepackage{xcolor}
\usepackage{amsmath,amssymb}
\usepackage{adjustbox}

%% Colors
\definecolor{rootcol}{HTML}{005F73}
\definecolor{c1}{HTML}{0077B6}
\definecolor{c2}{HTML}{E07000}
\definecolor{c3}{HTML}{AE2012}
\definecolor{c4}{HTML}{6A0572}
\definecolor{c5}{HTML}{1B7A3A}
\definecolor{c6}{HTML}{007F8C}

\forestset{
  branchstyle/.style 2 args={
    rounded corners=3pt, draw=#1, fill=#1!25, text=#1!70!black,
    font=\small\bfseries, align=center, inner sep=4pt,
    edge={draw=#1, thick},
  },
  leafstyle/.style 2 args={
    rounded corners=3pt, draw=#1!60, fill=#1!8, text=black,
    font=\scriptsize, align=left, inner sep=3pt,
    edge={draw=#1!60},
  },
}

\begin{document}
\pagestyle{empty}

\begin{adjustbox}{max width=\textwidth, max totalheight=\textheight, center}
\begin{forest}
  for tree={
    grow=east,
    l sep=7mm,
    s sep=2mm,
    edge={thick},
  },
  %% ROOT
  [{Cap.\ 13 --- Semi-Supervised Learning},
    rounded corners=4pt, draw=rootcol, fill=rootcol, text=white,
    font=\normalsize\bfseries, align=center, inner sep=6pt,
    %% BRANCH 1
    [{Background \& Assunzioni}, branchstyle={c1}{},
      [{Problema: pochi labeled $l$, molti unlabeled $u\gg l$}, leafstyle={c1}{}]
      [{Active Learning: il modello seleziona attivamente quali campioni etichettare}, leafstyle={c1}{}]
      [{Transductive: predice solo gli unlabeled visti in training}, leafstyle={c1}{}]
      [{Pure SSL: apprende una funzione generalizzabile a nuovi campioni}, leafstyle={c1}{}]
      [{Cluster Assumption: stesso cluster $\Rightarrow$ stessa classe}, leafstyle={c1}{}]
      [{Manifold Assumption: campioni vicini sulla varieta' $\Rightarrow$ output simili}, leafstyle={c1}{}]
    ]
    %% BRANCH 2
    [{Metodi Generativi (GMM)}, branchstyle={c2}{},
      [{Modello: $p(x)=\sum_i\alpha_i\mathcal{N}(x;\mu_i,\Sigma_i)$, un componente Gaussiano per classe}, leafstyle={c2}{}]
      [{Algoritmo EM su labeled $+$ unlabeled}, leafstyle={c2}{}]
      [{E-step: calcola $P(\text{componente}\mid x_j)$ per campioni non etichettati}, leafstyle={c2}{}]
      [{M-step: aggiorna $\alpha_i,\mu_i,\Sigma_i$ usando tutti i campioni}, leafstyle={c2}{}]
      [{Pro: semplice, efficace con pochissime etichette}, leafstyle={c2}{}]
      [{Contro: dipende fortemente dalla correttezza del modello generativo}, leafstyle={c2}{}]
    ]
    %% BRANCH 3
    [{Semi-Supervised SVM (S3VM)}, branchstyle={c3}{},
      [{Principio: low-density separation --- il confine attraversa zone a bassa densita'}, leafstyle={c3}{}]
      [{TSVM (Transductive SVM, Joachims 1999): assegna pseudo-label agli unlabeled}, leafstyle={c3}{}]
      [{Ricerca locale iterativa: scambia etichette di coppie probabilmente errate}, leafstyle={c3}{}]
      [{Aumenta gradualmente $C_u$ fino a $C_u=C_l$ per bilanciare l'influenza degli unlabeled}, leafstyle={c3}{}]
      [{Class imbalance: split $C_u$ in $C_u^+$ e $C_u^-$}, leafstyle={c3}{}]
      [{Varianti efficienti: LDS (gradient descent su graph kernel), meanS3VM}, leafstyle={c3}{}]
    ]
    %% BRANCH 4
    [{Graph-based SSL}, branchstyle={c4}{},
      [{Grafo $G=(V,E)$: nodi = campioni, archi pesati per similarita'}, leafstyle={c4}{}]
      [{Affinita': $W_{ij}=\exp(-\|x_i-x_j\|^2/2\sigma^2)$}, leafstyle={c4}{}]
      [{Label Propagation binaria (Zhu et al.\ 2003): energia $E(f)=\frac{1}{2}\sum_{ij}W_{ij}(f_i-f_j)^2$, sol.\ chiusa}, leafstyle={c4}{}]
      [{Label Propagation multiclasse (Zhou et al.\ 2004): $F^{(t+1)}=\alpha SF^{(t)}+(1-\alpha)Y$}, leafstyle={c4}{}]
      [{Converge a $F^*=(1-\alpha)(I-\alpha S)^{-1}Y$, minimizza rischio regolarizzato}, leafstyle={c4}{}]
      [{Limiti: memoria $O(n^2)$; difficile estendere a nuovi campioni}, leafstyle={c4}{}]
    ]
    %% BRANCH 5
    [{Disagreement-based Methods}, branchstyle={c5}{},
      [{Co-Training (Blum \& Mitchell 1998): due viste sufficienti e cond.\ indipendenti}, leafstyle={c5}{}]
      [{view1 $\to$ model1, view2 $\to$ model2; ognuno etichetta i campioni piu' confidenti per l'altro}, leafstyle={c5}{}]
      [{Teorema: performance boostata arbitrariamente con unlabeled (viste indipendenti)}, leafstyle={c5}{}]
      [{Varianti single-view: diversi algoritmi, diverso sampling, diversi parametri}, leafstyle={c5}{}]
      [{Condizione chiave: disaccordo significativo tra i learner (multi-view non necessario)}, leafstyle={c5}{}]
    ]
    %% BRANCH 6
    [{Semi-Supervised Clustering}, branchstyle={c6}{},
      [{Must-link: due campioni DEVONO stare nello stesso cluster}, leafstyle={c6}{}]
      [{Cannot-link: due campioni NON devono stare nello stesso cluster}, leafstyle={c6}{}]
      [{Constrained k-means (Wagstaff 2001): estende k-means rispettando i vincoli; violazione $\to$ cluster secondario}, leafstyle={c6}{}]
      [{Constrained Seed k-means (Basu 2002): labeled come semi per inizializzare i $k$ centroidi; semi fissi}, leafstyle={c6}{}]
    ]
  ]
\end{forest}
\end{adjustbox}

\end{document}
